\documentclass[11pt,a4paper]{article}
\usepackage[utf8]{inputenc}
\usepackage[french]{babel}
\usepackage[T1]{fontenc}
\usepackage{lmodern}
\usepackage[margin=1.5cm]{geometry}

\begin{document}

% En-tête
\begin{center}
    {\Huge \textbf{Lucas Bernard}} \\[0.2cm]
    {\Large \textbf{Ingénieur DevOps}} \\[0.1cm]
    lucas.bernard@email.com \ | \ 06 71 88 15 45
\end{center}

\vspace{0.3cm}

% Résumé / Profil
\section*{Profil Professionnel}
\noindent
Ingénieur DevOps et Cloud certifié, expert dans l'automatisation des infrastructures et la mise en place de pipelines CI/CD. J'aide les équipes de développement à livrer plus rapidement et en toute sécurité en adoptant les principes de l'Infrastructure as Code (IaC) et la conteneurisation des applications. Je milite pour la culture DevOps au quotidien.

\vspace{0.3cm}

% Compétences
\section*{Compétences Techniques}
\noindent
\textbf{GCP} $\bullet$ \textbf{Python} $\bullet$ \textbf{Docker} $\bullet$ \textbf{AWS} $\bullet$ \textbf{Linux} $\bullet$ \textbf{Bash} $\bullet$ \textbf{Prometheus} $\bullet$ \textbf{Kubernetes}

\vspace{0.3cm}

% Expériences
\section*{Expériences Professionnelles}
\noindent \textbf{Ingénieur DevOps / Confirmé} --- \textit{Datadog} \hfill \textbf{08/2019 à Aujourd'hui} \\
Mise en place de pipelines d'intégration et de déploiement continus (CI/CD) sur GitLab automatisant l'ensemble des tests et le déploiement sur les clusters Kubernetes. Astreinte technique (On-call) et résolution d'incidents critiques en production de niveau 3 (N3), avec rédaction de post-mortems détaillés (blameless). Implémentation d'un système de monitoring complet avec Prometheus et Grafana, réduisant le temps de détection des pannes (MTTD) de 60\%. Automatisation du provisionnement et de la configuration des serveurs Linux à l'aide de playbooks Ansible, éliminant les erreurs manuelles. \\[0.3cm]
\noindent \textbf{Ingénieur DevOps} --- \textit{Capgemini} \hfill \textbf{10/2018 à 05/2019} \\
Mise en place d'une architecture de journalisation centralisée (ELK Stack), facilitant l'analyse des logs et le diagnostic des anomalies en production. Gestion et sécurisation des secrets de l'entreprise via HashiCorp Vault, avec rotation automatique des clés d'accès des API tierces. Migration de 50 applications legacy vers une architecture conteneurisée sous Docker, réduisant drastiquement les coûts d'hébergement annuels. Audit de l'infrastructure cloud et application des principes du FinOps, permettant une réduction de la facture mensuelle AWS de 25\%. \\[0.3cm]
\noindent \textbf{Stage — Assistant Ingénieur DevOps} --- \textit{Capgemini} \hfill \textbf{01/2017 à 08/2018} \\
Déploiement d'une politique de sécurité 'Zero Trust', incluant la gestion stricte des identités (IAM) et le chiffrement des données au repos. Conception d'une infrastructure cloud résiliente sur AWS via Terraform (IaC), assurant une haute disponibilité des services et une scalabilité horizontale automatique. Mise en place d'une stratégie de sauvegarde automatisée (Disaster Recovery Plan) avec des tests de restauration trimestriels réussis à 100\%. \\[0.3cm]


\vspace{0.3cm}

% Formations
\section*{Formation \& Diplômes}
\noindent \textbf{Diplôme d'Ingénieur en Cloud Computing} \hfill \textbf{2015 - 2018} \\
\textit{École 42} \\[0.3cm]
\noindent \textbf{Classes Préparatoires (CPGE) en Filière Scientifique} \hfill \textbf{2013 - 2015} \\
\textit{Lycée Scientifique} \\[0.3cm]


\end{document}